\documentclass[11pt,a4paper]{article}
\usepackage[utf8]{inputenc}
\usepackage[margin=2cm]{geometry}
\usepackage{booktabs}
\usepackage{array}
\usepackage{amsmath}
\usepackage{amssymb}
\usepackage{graphicx}

% Define compact domain notation commands
\newcommand{\Realdomain}[1]{\ensuremath{[#1]\subset\mathbb{R}}}
\newcommand{\Intdomain}[1]{\ensuremath{[#1]\subset\mathbb{N}}}
\newcommand{\Catdomain}[1]{\ensuremath{\{\;}#1\ensuremath{\;\}}}

\title{Parameter Space for MOEADDouble Algorithm}
\author{jMetal Evolver Framework}
\date{\today}

\begin{document}

\maketitle

\section{Parameter Space}

\begin{table}[!tb]
\caption{Parameter space of MOEADDouble.}
\label{tab:params_moeaddouble}
\centering
\resizebox{\textwidth}{!}{%
\small
\begin{tabular}{r@{\hskip 1em}p{25em}}
\toprule
\bf Parameter & \bf Domain \\
\midrule
neighborhoodSize & \Intdomain{5, 50} \\
maximumNumberOfReplacedSolutions & \Intdomain{1, 5} \\
aggregationFunction & \Catdomain{tschebyscheff, weightedSum, penaltyBoundaryIntersection, modifiedTschebyscheff} \\
normalizeObjectives & \Catdomain{true, false} \\
epsilonParameterForNormalization & \Realdomain{1e-08, 25.0} \hspace{6em}\textbf{if} normalizeObjectives$=$true \\
pbiTheta & \Realdomain{1.0, 200.0} \hspace{6em}\textbf{if} aggregationFunction$=$penaltyBoundaryIntersection \\
algorithmResult & \Catdomain{externalArchive, population} \\
archiveType & \Catdomain{crowdingDistanceArchive, unboundedArchive} \hspace{6em}\textbf{if} algorithmResult$=$externalArchive \\
subProblemIdGenerator & \Catdomain{randomPermutationCycle, cyclicIntegerSequence} \\
createInitialSolutions & \Catdomain{default, latinHypercubeSampling, scatterSearch} \\
variation & \Catdomain{crossoverAndMutationVariation, differentialEvolutionVariation} \\
mutation & \Catdomain{uniform, polynomial, linkedPolynomial, nonUniform} \\
mutationProbabilityFactor & \Realdomain{0.0, 2.0} \\
mutationRepairStrategy & \Catdomain{random, round, bounds} \\
uniformMutationPerturbation & \Realdomain{0.0, 1.0} \hspace{6em}\textbf{if} mutation$=$uniform \\
polynomialMutationDistributionIndex & \Realdomain{5.0, 400.0} \hspace{6em}\textbf{if} mutation$=$polynomial \\
linkedPolynomialMutationDistributionIndex & \Realdomain{5.0, 400.0} \hspace{6em}\textbf{if} mutation$=$linkedPolynomial \\
nonUniformMutationPerturbation & \Realdomain{0.0, 1.0} \hspace{6em}\textbf{if} mutation$=$nonUniform \\
crossover & \Catdomain{SBX, blxAlpha, wholeArithmetic} \hspace{6em}\textbf{if} variation$=$crossoverAndMutationVariation \\
crossoverProbability & \Realdomain{0.0, 1.0} \hspace{6em}\textbf{if} variation$=$crossoverAndMutationVariation \\
crossoverRepairStrategy & \Catdomain{random, round, bounds} \hspace{6em}\textbf{if} variation$=$crossoverAndMutationVariation \\
sbxDistributionIndex & \Realdomain{5.0, 400.0} \hspace{6em}\textbf{if} crossover$=$SBX \\
blxAlphaCrossoverAlpha & \Realdomain{0.0, 1.0} \hspace{6em}\textbf{if} crossover$=$blxAlpha \\
blxAlphaBetaCrossoverBeta & \Realdomain{0.0, 1.0} \hspace{6em}\textbf{if} crossover$=$wholeArithmetic \\
blxAlphaBetaCrossoverAlpha & \Realdomain{0.0, 1.0} \hspace{6em}\textbf{if} crossover$=$wholeArithmetic \\
differentialEvolutionCrossover & \Catdomain{RAND_1_BIN, RAND_1_EXP, RAND_2_BIN} \hspace{6em}\textbf{if} variation$=$differentialEvolutionVariation \\
CR & \Realdomain{0.0, 1.0} \hspace{6em}\textbf{if} variation$=$differentialEvolutionVariation \\
F & \Realdomain{0.0, 1.0} \hspace{6em}\textbf{if} variation$=$differentialEvolutionVariation \\
selection & \Catdomain{populationAndNeighborhoodMatingPoolSelection} \\
neighborhoodSelectionProbability & \Realdomain{0.0, 1.0} \hspace{6em}\textbf{if} selection$=$populationAndNeighborhoodMatingPoolSelection \\
\bottomrule
\end{tabular}}
\end{table}

\section{Summary}

This parameter space contains 30 parameters total, of which 22 are conditional parameters that are only active under specific conditions.

\subsection{Domain Notation}

\begin{itemize}
\item \texttt{\textbackslash Realdomain{a, b}}: Real-valued parameters in range [a, b] $\subset \mathbb{R}$
\item \texttt{\textbackslash Intdomain{a, b}}: Integer-valued parameters in range [a, b] $\subset \mathbb{N}$
\item \texttt{\textbackslash Catdomain{opt1, opt2, ...}}: Categorical parameters with discrete options
\end{itemize}

\end{document}
